\chapter{Modelo de comportamiento}
\label{cap:comportamiento}

El modelo de comportamiento describe qué debe hacer el sistema internamente para responder a
los acontecimientos del modelo ambiental \parencite{yourdon1989}. Se compone de los diagramas
de flujo de datos, que muestran las funciones y los datos que intercambian; el diagrama
entidad-relación, que muestra la estructura de los datos almacenados; las mini
especificaciones, que describen la lógica de cada proceso elemental; y el diccionario de
datos, que define con precisión cada flujo, almacén y dato elemental.

% =========================================================
\section{Diagramas de flujo de datos}
\label{sec:dfd}

Los diagramas de flujo de datos (DFD) se construyeron con la notación de
\textcite{yourdon1989} y \textcite{demarco1979}:

\begin{itemize}
    \item Los \textbf{procesos} se dibujan como círculos, con un número jerárquico y un nombre
    de la forma verbo más objeto. El proceso 4.3 del nivel~1 se descompone en los procesos
    4.3.1, 4.3.2 y 4.3.3 del nivel~2, y así sucesivamente.
    \item Los \textbf{terminadores} se dibujan como rectángulos y los \textbf{almacenes} como
    dos líneas paralelas, identificados con la letra D y un número (D1 a D9).
    \item Los \textbf{flujos} son flechas rotuladas con el nombre del dato que transportan.
    Los flujos hacia o desde un almacén no se rotulan cuando transportan registros completos
    del almacén; la flecha de doble punta indica que el proceso lee y actualiza el almacén.
    \item Para evitar cruces, algunos terminadores y almacenes se dibujan más de una vez; las
    copias se marcan con asterisco (por ejemplo, \este{Jefe de faena*} o \este{D2* Clientes}).
    \item En los diagramas hijos, los flujos que vienen de otro proceso del diagrama padre o
    van hacia él se indican con el texto \este{desde el proceso n} o \este{hacia el proceso n}.
\end{itemize}

Cada diagrama hijo está \textbf{balanceado} con su padre: los flujos que entran y salen de un
proceso en el nivel superior son los mismos que entran y salen de su diagrama hijo, ya sea
directamente o mediante sus componentes. Por ejemplo, el flujo \este{solicitud de mantención}
del nivel~0 se descompone en el nivel~1 en \este{datos de usuario}, \este{datos de rol},
\este{datos de tipo de bulto}, \este{datos de cliente} y \este{datos de trabajador}, tal como
lo define el diccionario de datos (Sección~\ref{sec:diccionario}).

% ---------------------------------------------------------
\subsection{DFD de nivel 0}
\label{sec:dfd0}

El DFD de nivel~0 (Figura~\ref{fig:dfd0}) descompone el proceso único del diagrama de contexto
en los seis procesos principales del sistema y muestra los nueve almacenes de datos que
comparten. Se lee de izquierda a derecha siguiendo el ciclo de vida de una operación:

\begin{enumerate}
    \item \textbf{Administrar datos maestros} mantiene los almacenes de usuarios (D1), clientes
    (D2), trabajadores y cuadrillas (D3) y tipos de bulto (D4), que los demás procesos solo
    consultan.
    \item \textbf{Importar manifiesto} obtiene el manifiesto XML del Servicio Nacional de
    Aduanas y registra sus contenedores y líneas de carga en D5.
    \item \textbf{Programar faenas} crea la programación de cada contenedor en D6, asigna la
    cuadrilla y el jefe tarjador y entrega a este último sus faenas asignadas.
    \item \textbf{Ejecutar operación en terreno} registra el avance de la faena en D7 y la
    evidencia fotográfica en D8; al cerrar la faena emite el flujo interno \este{operación
    finalizada}.
    \item \textbf{Validar reporte} genera el borrador del reporte en D9, lo presenta al
    Supervisor y lo aprueba o rechaza; al aprobarlo emite el flujo interno \este{reporte
    aprobado}.
    \item \textbf{Distribuir y consultar reportes} envía el reporte aprobado al cliente,
    permite su reenvío y atiende las consultas del Jefe de faena y del Cliente del depósito.
\end{enumerate}

\diagramaapaisado{dfd_nivel0}
    {DFD de nivel 0}
    {Diagrama de flujo de datos de nivel 0 del sistema BeeTracer.}
    {fig:dfd0}

% ---------------------------------------------------------
\subsection{DFD de nivel 1}
\label{sec:dfd1}

En el nivel~1 se descompone cada proceso del nivel~0 en los subprocesos que lo componen. A continuación se describe cada diagrama; los seis diagramas se presentan juntos al final de esta sección (Figuras~\ref{fig:dfd1} a~\ref{fig:dfd6}).

\paragraph{Proceso 1: Administrar datos maestros.} La solicitud de mantención del
Administrador se separa según el tipo de dato maestro (Figura~\ref{fig:dfd1}). Cada
subproceso valida los datos, aplica la acción solicitada (crear, editar, eliminar o cambiar
el estado) sobre su almacén y devuelve el resultado. El proceso 1.1 consulta además los
almacenes de clientes y trabajadores, porque una cuenta de usuario puede quedar ligada a la
ficha de un trabajador (jefe tarjador) o de un cliente.

\paragraph{Proceso 2: Importar manifiesto.} El proceso 2.1 verifica que el viaje no haya sido
importado antes y solicita el manifiesto a Aduanas; el 2.2 valida la estructura y los datos
del XML recibido; si es válido, el 2.3 lo registra y, si no, el 2.4 genera el reporte de
errores sin registrar información parcial. En ambos casos el 2.5 notifica el resultado a
Aduanas y entrega el resumen al Jefe de faena (Figura~\ref{fig:dfd2}).

\paragraph{Proceso 3: Programar faenas.} El Jefe de faena crea la programación de un
contenedor pendiente (3.1), le asigna la cuadrilla y el jefe tarjador (3.2) y puede
reprogramarla, anularla, eliminarla o cambiar su estado (3.3). Cuando el jefe tarjador inicia
sesión en la aplicación móvil, el proceso 3.4 le entrega las faenas que tiene asignadas, con
los datos del contenedor y su carga declarada (Figura~\ref{fig:dfd3}).

\paragraph{Proceso 4: Ejecutar operación en terreno.} Cada etapa del flujo guiado del jefe
tarjador corresponde a un subproceso (Figura~\ref{fig:dfd4}): iniciar la operación (4.1),
verificar el sello (4.2), registrar cada lote separado por cliente (4.3), registrar las
incidencias (4.4) y cerrar la faena (4.5). Los procesos que reciben una fotografía la
registran en el almacén D8 junto con su etapa, fecha, hora y geoposición, lo que materializa
la regla de evidencia obligatoria. Al cerrar la faena, el proceso 4.5 emite el flujo
\este{operación finalizada} hacia el proceso 5.

\paragraph{Proceso 5: Validar reporte.} El proceso 5.1 genera el borrador del reporte cuando
recibe una operación finalizada o una evidencia corregida. El Supervisor solicita revisar un
reporte (5.2) y lo aprueba (5.3) o rechaza la fotografía defectuosa (5.4), lo que genera una
solicitud de recaptura al jefe tarjador. La fotografía de reemplazo se registra en el proceso
5.5, que reactiva la generación de una nueva versión del reporte (Figura~\ref{fig:dfd5}).

\paragraph{Proceso 6: Distribuir y consultar reportes.} El proceso 6.1 envía el reporte
aprobado a los correos del cliente, ya sea automáticamente al aprobarse o a solicitud del Jefe
de faena, y registra cada envío. El Jefe de faena consulta las operaciones y modifica su
estado (6.2), y el Cliente del depósito consulta el estado de sus operaciones (6.3) y
descarga el reporte aprobado (6.4), viendo solo las operaciones con carga consignada a su
empresa (Figura~\ref{fig:dfd6}).

\diagramaapaisado{dfd1_maestros}
    {DFD de nivel 1 - Proceso 1}
    {DFD de nivel 1 - Proceso 1: Administrar datos maestros.}
    {fig:dfd1}

\diagramaapaisado{dfd2_importar}
    {DFD de nivel 1 - Proceso 2}
    {DFD de nivel 1 - Proceso 2: Importar manifiesto.}
    {fig:dfd2}

\diagramaapaisado{dfd3_programar}
    {DFD de nivel 1 - Proceso 3}
    {DFD de nivel 1 - Proceso 3: Programar faenas.}
    {fig:dfd3}

\diagramaapaisado{dfd4_operacion}
    {DFD de nivel 1 - Proceso 4}
    {DFD de nivel 1 - Proceso 4: Ejecutar operación en terreno.}
    {fig:dfd4}

\diagramaapaisado{dfd5_validar}
    {DFD de nivel 1 - Proceso 5}
    {DFD de nivel 1 - Proceso 5: Validar reporte.}
    {fig:dfd5}

\diagramaapaisado{dfd6_distribuir}
    {DFD de nivel 1 - Proceso 6}
    {DFD de nivel 1 - Proceso 6: Distribuir y consultar reportes.}
    {fig:dfd6}

% ---------------------------------------------------------
\subsection{DFD de nivel 2}
\label{sec:dfd2}

Se descomponen en un nivel más los tres procesos de nivel~1 cuya lógica involucra varios
pasos con almacenes distintos. Los demás procesos de nivel~1 se consideran
\emph{elementales}: su lógica cabe en una mini especificación de una página.

\paragraph{Proceso 2.3: Registrar manifiesto.} El registro del manifiesto válido se realiza en
cadena (Figura~\ref{fig:dfd23}): primero la cabecera (nave, viaje y número), luego cada
contenedor declarado y, por cada línea de carga, se verifica que el consignatario exista en
el almacén de clientes, creándolo con sus datos mínimos (RUT y razón social) si no existe,
antes de registrar la línea con su tipo de bulto.

\diagramatikz{dfd23_registrar_manifiesto}
    {DFD de nivel 2 - Proceso 2.3}
    {DFD de nivel 2 - Proceso 2.3: Registrar manifiesto.}
    {fig:dfd23}

\paragraph{Proceso 3.2: Asignar cuadrilla y jefe tarjador.} La asignación es un diálogo con el
Jefe de faena (Figura~\ref{fig:dfd32}): el sistema determina qué cuadrillas y jefes tarjadores
están habilitados y libres en la franja de la programación (3.2.1); el Jefe de faena elige y
el sistema vuelve a verificar la disponibilidad, porque otra programación pudo ocuparlos
entretanto (3.2.2); finalmente registra la asignación (3.2.3).

\diagramatikz{dfd32_asignar}
    {DFD de nivel 2 - Proceso 3.2}
    {DFD de nivel 2 - Proceso 3.2: Asignar cuadrilla y jefe tarjador.}
    {fig:dfd32}

\paragraph{Proceso 4.3: Registrar lote.} Por cada lote separado (Figura~\ref{fig:dfd43}), el
sistema identifica la línea de carga a la que pertenece o, en tarja ciega, la crea con los
datos ingresados (4.3.1); contrasta la cantidad recibida con la declarada y registra la
fotografía del lote (4.3.2); y evalúa la diferencia para confirmar el lote o exigir el
registro de una incidencia de faltante, sobrante o daño (4.3.3).

\diagramatikz{dfd43_lote}
    {DFD de nivel 2 - Proceso 4.3}
    {DFD de nivel 2 - Proceso 4.3: Registrar lote.}
    {fig:dfd43}
